\documentclass[11pt]{article}
\usepackage[review]{acl}
\usepackage{times}
\usepackage{latexsym}
\usepackage[T1]{fontenc}
\usepackage[utf8]{inputenc}
\usepackage{microtype}
\usepackage{inconsolata}
\usepackage{lipsum}

\title{A Multi-Pronged Approach to Detecting Hallucinated References in Academic Papers}

\author{First Author \\
  Test University \\
  Department of Computer Science \\
  \texttt{first@test.edu} \\\And
  Second Author \\
  Test University \\
  Department of Computer Science \\
  \texttt{second@test.edu} \\}

\begin{document}
\maketitle

\begin{abstract}
We present a comprehensive system for detecting hallucinated references in academic papers. Our approach combines multiple PDF extraction backends with a multi-API verification pipeline and agentic retry strategies. We evaluate on a dataset of papers containing both real and fabricated citations, achieving high precision and recall. The system produces a ranked list of references by hallucination likelihood, enabling efficient review of potentially problematic citations. Our tool is fully open source and processes all documents locally for privacy.
\end{abstract}

\section{Introduction}

Natural language processing has seen remarkable advances in recent years, driven by the development of large pretrained language models \cite{devlin-etal-2019-bert,vaswani-etal-2017-attention,brown-etal-2020-language}. These models have demonstrated impressive capabilities across a wide range of tasks, from text generation to question answering and beyond.

\lipsum[1-3]

The problem of hallucinated references is particularly concerning in academic publishing. Language models can generate plausible-looking citations that combine real author names with fictitious titles, or produce entirely fabricated references that pass cursory human review \cite{liu2019roberta,raffel-etal-2020-exploring}.

\lipsum[4-5]

\section{Background}

\subsection{Reference Extraction from Scientific Documents}

Extracting structured references from PDF documents has been studied extensively. Various tools exist for this purpose, ranging from rule-based systems to deep learning approaches \cite{wolf-etal-2020-transformers}. The accuracy of reference extraction directly impacts downstream verification tasks.

\lipsum[6-8]

\subsection{Citation Verification Methods}

Prior work on citation verification has focused primarily on detecting broken URLs and missing DOIs. More recent approaches leverage academic search APIs to verify that cited papers actually exist \cite{lewis-etal-2020-bart,mikolov-etal-2013-distributed}.

\lipsum[9-11]

\subsection{Hallucination Detection in Language Model Outputs}

Large language models are known to produce hallucinated content, including fabricated citations \cite{pennington-etal-2014-glove}. Studies have shown that GPT-based models can generate plausible-looking references that combine real author names with non-existent paper titles.

\lipsum[12-14]

\section{Methodology}

\subsection{Multi-Pronged Reference Extraction}

Our approach uses multiple extraction backends operating in parallel. Each backend has different strengths for handling various PDF formatting challenges \cite{peters-etal-2018-deep}.

\lipsum[15-17]

\subsection{Field Parsing Pipeline}

Individual reference strings are parsed into structured fields using a cascade of parsers, from deterministic regex patterns to flexible LLM-based parsing \cite{devlin-etal-2019-bert,vaswani-etal-2017-attention}.

\lipsum[18-20]

\subsection{Multi-API Verification}

Extracted references are verified against multiple academic databases including Semantic Scholar, CrossRef, DBLP, OpenAlex, and the ACL Anthology \cite{brown-etal-2020-language,liu2019roberta}.

\lipsum[21-23]

\subsection{Hallucination Scoring}

We compute a composite hallucination score based on weighted signals including title similarity, author overlap, author order correctness, year matching, and cross-API consensus \cite{raffel-etal-2020-exploring}.

\lipsum[24-26]

\subsection{Agentic Retry Strategies}

When initial API searches fail to find a match, our system employs a series of query reformulation strategies \cite{lewis-etal-2020-bart,mikolov-etal-2013-distributed}.

\lipsum[42-44]

\section{Experiments}

We evaluated our system on a dataset of 100 papers containing a mix of real and hallucinated references \cite{wolf-etal-2020-transformers,lewis-etal-2020-bart}.

\lipsum[26-28]

\subsection{Results}

Our system achieved strong performance across all evaluation metrics. The multi-API approach proved essential for reliable verification \cite{mikolov-etal-2013-distributed,pennington-etal-2014-glove}.

\lipsum[29-31]

\subsection{Analysis}

The agentic retry strategies significantly improved recall compared to direct title matching alone \cite{peters-etal-2018-deep}.

\lipsum[32-33]

\subsection{Ablation Study}

We conducted an ablation study to evaluate the contribution of each component to overall system performance.

\lipsum[45-46]

\section{Discussion}

\lipsum[34-36]

\subsection{Limitations}

\lipsum[37-38]

\subsection{Ethical Considerations}

\lipsum[47-48]

\section{Related Work}

Several tools have been developed for reference verification, each with different trade-offs between accuracy and usability \cite{devlin-etal-2019-bert,brown-etal-2020-language,raffel-etal-2020-exploring}.

\lipsum[39-40]

\lipsum[49-50]

\section{Conclusion}

We presented a comprehensive system for detecting hallucinated references in academic papers \cite{vaswani-etal-2017-attention,wolf-etal-2020-transformers}. Our multi-pronged approach achieves high precision and recall while keeping all PDF processing fully local.

\lipsum[41]

\bibliography{test_refs}
\bibliographystyle{acl_natbib}

\end{document}
